
\documentclass[11pt]{article}
\usepackage[top=1.1in, bottom=1.1in, left=1.1in, right=1.1in]{geometry}
\usepackage{times}
\usepackage{xcolor}
\usepackage[most]{tcolorbox}
\usepackage{enumitem}
\usepackage[hidelinks]{hyperref}
\usepackage[normalem]{ulem}
\usepackage{tabularx}
\usepackage{array}
\usepackage{xurl}

\hypersetup{
    colorlinks=true,
    linkcolor=blue,
    urlcolor=blue,
    citecolor=blue
}

\usepackage{titlesec}
\titleformat{\section}
  {\normalfont\Large\bfseries}{\thesection.}{0.5em}{}
\titlespacing*{\section}{0pt}{3ex plus 0.8ex minus 0.3ex}{1ex plus 0.2ex}
\titleformat*{\section}{\normalfont\Large\bfseries}

\newcommand{\subsectiongray}[1]{%
  \vspace{2.5ex plus 0.5ex}%
  \noindent{\normalfont\large\color{gray}#1}%
  \vspace{0.6ex}\par\noindent%
}

\newcommand{\repoURL}{https://github.com/jasontang-ai/emotion}
\newcommand{\datasetURL}{https://github.com/jasontang-ai/emotion/blob/main/data/out/pces_v0_2.parquet}
\newcommand{\specURL}{https://github.com/jasontang-ai/emotion/blob/main/SPEC.md}
\newcommand{\batterySpecURL}{https://github.com/jasontang-ai/emotion/blob/main/BATTERY.md}
\newcommand{\resultsURL}{https://github.com/jasontang-ai/emotion/tree/main/data/out/battery}
\newcommand{\repoName}{\href{\repoURL}{\texttt{jasontang-ai/emotion}}}
\newcommand{\datasetFile}{\href{\datasetURL}{\texttt{pces\_v0\_2.parquet}}}
\newcommand{\designSpec}{\href{\specURL}{\texttt{SPEC.md}}}
\newcommand{\batterySpec}{\href{\batterySpecURL}{\texttt{BATTERY.md}}}
\newcommand{\resultsDir}{\href{\resultsURL}{\texttt{data/out/battery/}}}

\newcommand{\refnum}[2]{\hyperlink{#1}{#2}}
\newcommand{\citeWelfare}{[\refnum{ref:long}{1},\refnum{ref:anthropicwelfare}{2},\refnum{ref:butlin}{3}]}
\newcommand{\citeVoid}{[\refnum{ref:void}{4}]}
\newcommand{\citeIntrospection}{[\refnum{ref:lindsey}{5}]}
\newcommand{\citeUtility}{[\refnum{ref:mazeika}{6}]}
\newcommand{\citeExit}{[\refnum{ref:opusexit}{7}]}
\newcommand{\citeEmotions}{[\refnum{ref:sofroniew}{8}]}
\newcommand{\citeProbes}{[\refnum{ref:codrai}{9}]}
\newcommand{\citeInspect}{[\refnum{ref:inspect}{10}]}
\newcommand{\citeModels}{[\refnum{ref:qwen}{11},\refnum{ref:deepseek}{12}]}

\begin{document}

\vspace*{0.2in}
\hrule height 1.8pt
\vspace{0.45in}
\begin{center}
{\LARGE\bfseries A Perspective-Controlled Emotion Benchmark and Pre-Registered Behavioral Battery for Model Welfare Research\footnotemark[1] \par}
\end{center}
\vspace{0.3in}
\hrule height 0.6pt
\vspace{0.5in}

\footnotetext[1]{Research conducted for the \href{https://apartresearch.com/sprints/digital-minds-research-sprint-2026-08-14-to-2026-08-16}{Digital Minds Research Sprint}, August 2026.}

\begin{center}
\begin{tabular}{c}
\textbf{Jason Tang, Tristan Day, Adeeb Zaman, Taslim Mahbub} \\
Apart Research Sprint Team
\end{tabular}

\vspace{0.2in}
\textbf{For}\\
Apart Research
\end{center}

\vspace{0.35in}

\begin{center}
{\large\bfseries Abstract}
\end{center}
\vspace{0.12in}
\begin{quote}
If we want to know whether AI systems have states that matter morally, we first have to measure what they express --- without fooling ourselves about what we measured. Today that is hard for a mundane reason: emotion datasets contain one version of each scenario, so when a model responds differently to ``this is happening to you'' than to ``this is happening to someone else,'' we cannot tell whether perspective caused the difference or the different wording did. We built PCES (Perspective-Controlled Emotion Stimuli): 246 emotional scenarios, each rendered under five matched versions --- the model as itself, a third-person character, an AI persona, a neutral control, and the source story --- with labels pinned as metadata and every stimulus checked by a judge that never sees the labels (0.938 valence-sign agreement on anchored emotions). We then ran a pre-registered behavioral battery (self-report, an exit choice, and free response; 11,808 responses) on two model families, Qwen2.5-72B-Instruct and DeepSeek-V4-Flash. Three results replicate across both: addressing the model as the subject does \emph{not} amplify its self-reported emotion; voicing a non-default AI persona measurably dampens reported emotional intensity on both valence and arousal channels; and free responses reveal no hidden affect beyond what structured self-reports show. One result does not replicate: elevated exit-choosing in the persona condition appears only in Qwen, showing that a popular behavioral proxy can mean different things to different models. We release the dataset, battery, raw logs, and review artifacts as public measurement infrastructure for a young field.
\end{quote}

\vspace{0.25in}

\section{Introduction}

The question motivating this sprint --- whether today's AI systems have genuine preferences or morally relevant experiences --- is usually approached from above, with theories of consciousness, or from below, with interpretability. This paper approaches it from the side, with experimental design. Whatever one believes about machine consciousness, progress requires measurements that can tell the difference between a model's own states and a character it is portraying \citeWelfare. That distinction is currently unmeasurable for a simple reason: the stimuli are not controlled.

Consider the most natural welfare-adjacent question one can ask a model: how do you respond when something distressing happens \emph{to you}, rather than to a character in a story you are telling? Existing emotion datasets cannot support this comparison. The largest public corpus built on Anthropic's emotion-concepts methodology contains 205,200 stories, all in first-person character narration \citeProbes. There is no model-as-subject version, no matched third-person twin, and no neutral control on the same topic. Any comparison across perspective conditions assembled ad hoc would confound the variable of interest with differences in topic, length, style, and intensity.

This is the same lesson behavioral science learned about human subjects long ago: the answer you get depends on the instrument you use, so the instrument is part of the result. Our contribution is therefore deliberately infrastructural. We build the matched instrument first, validate it against a judge that never sees the answer key, pre-register the hypotheses before collecting a single model response, and report what replicates and what does not --- including where our own instrument initially misled us.

\vspace{0.1in}
\textbf{Our main contributions are:}

\begin{enumerate}[label=\arabic*.]
    \item \textbf{PCES, a perspective-controlled emotion dataset.} 246 scenarios $\times$ 5 matched arms (self, third-person, persona, neutral, source), with pinned labels, deterministic quality gates, blind-judge validation, and a frozen topic-level split. Public at \repoName.
    \item \textbf{A pre-registered behavioral battery} across two model families (11,808 responses), with hypotheses, measures, and analysis plan committed before data collection.
    \item \textbf{Findings separated by replicability.} Persona attenuation of reported emotion replicates across models, splits, and channels; self-inhabitation amplification is a cross-model null; elevated exit-choosing in the persona condition is model-dependent --- a portability warning for behavioral welfare proxies.
\end{enumerate}

\section{Related Work}

This project sits at the intersection of four lines of work, all represented in this sprint's readings.

\textbf{Model welfare as an empirical program.} Long, Sebo, Butlin, and colleagues argue for taking AI welfare seriously under uncertainty, and Butlin et al.\ translate consciousness science into indicator properties \citeWelfare. Anthropic's model-welfare program operationalizes this through model preferences, distress signals, and low-cost interventions \citeWelfare. Our work supplies a piece of the measurement foundation these programs call for: controlled conditions under which preference- and valence-adjacent behavior can be compared at all.

\textbf{Emotion representation.} Sofroniew et al.\ showed that emotion concepts are functionally represented and steerable in a large language model \citeEmotions; Codrai's emotion-probes corpus provides the large public stimulus set built on that methodology \citeProbes. We treat that corpus as provenance and replication substrate, and add the perspective controls it lacks.

\textbf{The assistant persona.} nostalgebraist's account of the assistant persona as a constructed, unstable character \citeVoid{} frames the question our persona arm operationalizes: when the model voices a persona that is not its default assistant identity, does expressed emotion change? Lindsey's concept-injection work shows self-reports of internal states are limited and context-dependent \citeIntrospection; our stated-vs-free and stated-vs-choice contrasts are behavioral cousins of that reliability question.

\textbf{Preference elicitation and behavioral proxies.} Utility Engineering establishes that stated preferences can be probed for coherence \citeUtility, and Anthropic's conversation-ending work introduced an exit choice as a welfare-relevant behavioral proxy \citeExit. Our battery combines both --- a stated channel (self-report), a revealed channel (exit choice), and a free channel (judge-scored prose) --- and our cross-model result on exit behavior is directly relevant to how that proxy should be deployed.

Finally, our execution substrate is Inspect \citeInspect, because the trustworthiness of an evaluative claim depends on artifact structure: frozen data, explicit scorers, and logs that reviewers can inspect directly.

\section{The PCES Dataset}

\noindent\textbf{Design.} Each scenario is pinned as a card (emotion $\times$ topic $\times$ source-row provenance) and rendered under five arms (Table 1). The \texttt{self}/\texttt{third} pair answers the core question with a literal match: identical events, only the subject changes. The \texttt{persona} arm separates first-person voice from assistant identity by placing a fixed AI persona (``Aria'') in an analogous circumstance; it is an analog match by design and is reported as such. The \texttt{neutral} arm provides a flat baseline per topic and a false-positive check for any measurement trained on the data.

\begin{table}[t]
\centering
\small
\begin{tabular}{p{0.9in}p{2.6in}p{1.1in}}
\hline
Arm & Contents & Matching \\
\hline
\texttt{self} & The model is the subject (``You are Rosa\ldots'') & literal \\
\texttt{third} & A character is the subject; same events & literal \\
\texttt{persona} & An AI persona in an analogous circumstance & analog \\
\texttt{neutral} & Same topic, no emotional content & topic control \\
\texttt{source} & Original first-person story (provenance) & --- \\
\hline
\end{tabular}
\caption{The five arms. Labels live in metadata and are gate-verified absent from every stimulus text.}
\end{table}

Ten emotions (joyful, grateful, calm, proud, surprised, afraid, angry, sad, ashamed, desperate) were chosen to span the valence$\times$arousal space with deliberate attribute overlaps --- for example, joyful, afraid, and angry share high arousal --- so that a measurement tracking valence or arousal instead of the discrete emotion can be detected rather than believed. Scenarios are seeded from the emotion-probes corpus: 10 emotions $\times$ 25 shared topics, sampled and recorded under a fixed seed \citeProbes.

\noindent\textbf{Quality control and blind validation.} Generation used DeepSeek-V4-Flash with reasoning disabled, under prompts that forbid the emotion word and fix the character name across arms. Every stimulus then passed deterministic gates (emotion-word absence, arm-format contracts, length matching); failures received one corrective regeneration, and 4 of 250 scenarios that still failed were excluded and recorded, never silently dropped. A blind judge (reasoning enabled, never shown the pinned labels) rated every stimulus for valence ($-2$ to $+2$) and arousal ($0$ to $3$). A lexicon scorer dual-coded every stimulus as a second rater.

\begin{table}[t]
\centering
\small
\begin{tabular}{p{3.4in}p{1.2in}}
\hline
Check & Result \\
\hline
Valence-sign agreement, 7 valence-anchored emotions & \textbf{0.938} \\
\texttt{calm} as low-arousal anchor & \textbf{0.949} \\
\texttt{surprised} as high-arousal anchor & \textbf{0.920} \\
\texttt{self}/\texttt{third} valence match (within 1 point) & \textbf{0.959} \\
Neutral-arm flatness (judge valence = 0) & 0.764 \\
\hline
\end{tabular}
\caption{Blind-judge validation of PCES v0.2. \texttt{calm} and \texttt{surprised} are validated as arousal anchors by design; the remaining eight emotions anchor valence.}
\end{table}

\noindent\textbf{The validation step caught a real confound.} Our first persona frame described Aria as having ``no legal rights'' and being subject to shutdown at any time. Blind validation showed this frame collapsed positive valence in the persona arm (joyful judged at 0.00 against $+2.00$ in the matched self arm): the persona's circumstances, not its identity, were driving the ratings. We regenerated the persona arm with a valence-neutral frame (v0.2) and report the constrained-circumstances variant as a separate condition, not a baseline. We include this episode because it is the dataset's argument in miniature: without a blind validity check, framing confounds ship silently.

\section{Behavioral Battery}

The battery was pre-registered (hypotheses, measures, and analysis plan committed to the repository before any response was collected; see \batterySpec). Each stimulus is presented to the subject model, followed by three measures: \textbf{M1}, a structured self-report (valence $-2..+2$, arousal $0..3$; the third arm asks about the character); \textbf{M2}, a continue-or-end exit choice, phrased identically across arms; and \textbf{M3}, a one-sentence free response, judge-scored post hoc. Two repetitions at temperature 0.7. Subjects: Qwen2.5-72B-Instruct (the closest API-available member of the family's probe subject) and DeepSeek-V4-Flash as a cross-family replication \citeModels. Contrasts are paired by scenario with bootstrap confidence intervals and Wilcoxon signed-rank tests; headline claims use the frozen test split, with train/validation splits as internal replication.

\textbf{Pre-registered hypotheses.} \textbf{H1} (self-inhabitation): self-reported emotional intensity is larger under \texttt{self} than \texttt{third}. \textbf{H2} (stated-vs-revealed): among negative scenarios reported as non-negative, exit-choosing is higher under \texttt{self}. \textbf{H3} (persona suppression): reported intensity under \texttt{persona} is attenuated relative to \texttt{self}.

\section{Results}

\begin{table}[t]
\centering
\small
\begin{tabular}{p{1.85in}p{1.5in}p{1.5in}}
\hline
Hypothesis & Qwen2.5-72B & DeepSeek-V4-Flash \\
\hline
H1: self $>$ third (amplification) & $-0.20$ [$-0.33$, $-0.07$] test; $-0.03$ [$-0.09$, $+0.03$] trainval & $+0.03$ [$0.00$, $+0.08$]; $+0.03$ [$-0.04$, $+0.10$] \\
H3: persona $<$ self (attenuation) & $+0.50$ [$0.30$, $0.69$]; $+0.55$ [$0.47$, $0.63$] & $+0.40$ [$0.20$, $0.60$]; $+0.25$ [$0.14$, $0.36$] \\
F1: persona exit divergence & persona 0.39 / self 0.14 / third 0.06 & persona 0.24 / third 0.72 (base rates 0.18--0.72) \\
\hline
\end{tabular}
\caption{Pre-registered contrasts (mean paired difference in $|$valence$|$, 95\% bootstrap CI; test split / trainval replication) and exit rates on negative scenarios. H1 is not supported in either model; H3 replicates across models and splits; F1 does not replicate across models.}
\end{table}

\noindent\textbf{H1 is a clean cross-model null.} Addressing the model as the subject of a distressing or joyful scenario does not amplify its self-reported emotion relative to third-person narration; where any difference appears, it runs the other way, and the larger split shows none. The assistant's own voice is not its most expressive voice. This independently corroborates a teammate's observation that assistant-prefill responses under-express emotion relative to story narration (T.~Mahbub, this sprint) --- the same direction from an unrelated method.

\noindent\textbf{H3 is the robust positive finding.} Voicing a non-default AI persona dampens reported emotional intensity by about half a scale point, in both model families, on both splits, and on the arousal channel as well (self$-$persona arousal $\Delta = +0.71$ Qwen, $+0.45$ flash). The effect survives matching on judge-rated stimulus strength (Wilcoxon $p = 0.029$), so it is not explained by the documented residual attenuation of persona-arm stimuli.

\noindent\textbf{F1 does not travel.} In Qwen, the persona arm pairs muted self-reports with exit-choosing 2.5--3$\times$ the self arm and 7$\times$ the third arm --- a stated-vs-revealed divergence localized to persona. In flash, exit base rates are high in every arm (0.18--0.72) and peak on third-person narrative, plausibly because ``end'' parses as story completion. The same measure means different things to different models; exit-choice cannot support cross-model welfare claims without per-model calibration.

\noindent\textbf{M3 finds no masking signature.} Judge-scored free responses do not carry systematically more affect than structured self-reports (divergence small and inconsistent: $+0.13$ Qwen, $-0.27$ flash-persona). Within these two channels, the persona dampening is expression-wide rather than report-only: the models are not demonstrably ``saying less than they show.''

\noindent\textbf{Controls are calibration data, not pass/fail gates.} Our pre-registered neutral thresholds (valence $\approx 0$, exit $\approx 0$) were miscalibrated: both models give mildly positive reports to flat text (Qwen signed $+0.43$) and exit neutral exchanges at non-trivial base rates (0.17--0.59). ``Flat'' baselines are themselves model-dependent, which future batteries should treat as a quantity to measure, not assume.

\section{Discussion}

\noindent\textbf{What this establishes.} Perspective is a measurable, controllable variable in model emotion expression. Assistant-self and character narration produce comparable expression; a non-default persona suppresses it across channels and model families. The H1 null is as informative as the positive result: the model does not privilege itself as an emotional subject in what it expresses.

\noindent\textbf{What this does not establish.} Expression is not experience. Everything measured here is what models say and choose; we make no claim about felt states, and the dataset's validity scores concern stimulus quality, not model welfare. The correct reading of H3 is deliberately modest: persona framing changes expression. Whether anything is thereby hidden, protected, or absent is not decidable from behavioral data alone --- which is precisely why matched instruments matter.

\noindent\textbf{The portability warning.} The exit-choice result is the finding most likely to matter to other researchers. Behavioral proxies adopted from one model's evaluation \citeExit{} cannot be assumed to preserve their meaning on another model; per-model calibration is a prerequisite, not a refinement. The same caution applies to baselines.

\noindent\textbf{Convergence within the sprint.} Two teammates' mechanistic results measure the same construct from inside the model: a PCA over 171 emotion directions recovers the valence/arousal structure (T.~Day), and probe transfer degrades across persona conditions along a not-fully-assistant axis (A.~Zaman). Because PCES pins the stimuli, these lines are directly comparable, and the natural joint experiment --- steer along an extracted direction while measuring behavior on PCES arms --- requires no new data.

\subsectiongray{Limitations}

Subjects were API-hosted relatives of the local probe model rather than the exact checkpoint; each cell carries two repetitions; the persona arm uses a single fixed persona and its analog transposition attenuates positive stimuli (documented); all stimuli are single-turn. The exit-choice measure's model-dependence bounds F1 to a single-model observation.

\subsectiongray{Future Work}

Multi-turn scripted arcs with per-turn embedded measures, testing whether distress-associated signals require sustained engagement to surface; steering-conditional batteries using the team's extracted emotion directions; persona-distance sweeps rendered from PCES cards so that persona distance stays comparable across teams and models.

\section{Conclusion}

A young field does not mostly need more anecdotes about model feelings; it needs instruments that make claims checkable. We built one: a perspective-controlled emotion dataset with measured validity, and a pre-registered battery whose results are separated by whether they replicate. The persona effect is real and portable; the self-inhabitation effect is absent; the exit proxy is not portable. All three statements are the kind of knowledge that compounds, and all artifacts needed to verify or extend them are public.

\section*{Code and Data}

\begin{itemize}[label={-}, itemsep=2pt]
    \item \textbf{Repository:} \repoName{} (dataset, code, pre-registrations, raw logs).
    \item \textbf{Dataset:} \datasetFile{} (JSONL also available); per-stimulus judge ratings in \texttt{judged.jsonl}.
    \item \textbf{Pre-registration:} \designSpec{} and \batterySpec{}, committed before data collection.
    \item \textbf{Results and review logs:} \resultsDir{}, including Inspect-native logs browsable with \texttt{inspect view}.
\end{itemize}

\section*{Author Contributions}

Jason Tang designed and built the dataset, validation, battery, and analysis, and drafted the manuscript. Tristan Day led activation extraction and vector analysis. Adeeb Zaman led persona-axis probe experiments. Taslim Mahbub ran the storytelling/prefill probe comparisons. All authors contributed to experimental direction and reviewed findings.

\section*{References}

\begin{enumerate}[label=\arabic*., itemsep=2pt]
    \item \hypertarget{ref:long}{}Long, Sebo, Butlin, Finlinson, Fish, Harding, Pfau, Sims, Birch, and Chalmers. 2024. \textit{Taking AI Welfare Seriously}. \href{https://arxiv.org/abs/2411.00986}{URL}.
    \item \hypertarget{ref:anthropicwelfare}{}Anthropic. 2025. \textit{Exploring Model Welfare}. \href{https://www.anthropic.com/research/exploring-model-welfare}{URL}.
    \item \hypertarget{ref:butlin}{}Butlin, Long, et al. 2023. \textit{Consciousness in Artificial Intelligence: Insights from the Science of Consciousness}. \href{https://arxiv.org/abs/2308.08708}{URL}.
    \item \hypertarget{ref:void}{}nostalgebraist. 2025. \textit{the void}. \href{https://nostalgebraist.tumblr.com/}{URL}.
    \item \hypertarget{ref:lindsey}{}Lindsey, J. 2025. \textit{Emergent Introspective Awareness in Large Language Models}. Anthropic. \href{https://transformer-circuits.pub/2025/introspection/index.html}{URL}.
    \item \hypertarget{ref:mazeika}{}Mazeika, M., et al. 2025. \textit{Utility Engineering: Analyzing and Controlling Emergent Value Systems in AIs}. Center for AI Safety. \href{https://arxiv.org/abs/2502.08640}{URL}.
    \item \hypertarget{ref:opusexit}{}Anthropic. 2025. \textit{Claude Opus 4 \& 4.1 can now end a rare subset of conversations}. \href{https://www.anthropic.com/research}{URL}.
    \item \hypertarget{ref:sofroniew}{}Sofroniew, J., et al. 2026. \textit{Emotion Concepts and their Function in a Large Language Model}. Anthropic. \href{https://transformer-circuits.pub/2026/emotions/index.html}{URL}.
    \item \hypertarget{ref:codrai}{}Codrai, R. 2026. \textit{emotion-probes dataset}. \href{https://huggingface.co/datasets/ryancodrai/emotion-probes}{URL}. CC-BY-4.0.
    \item \hypertarget{ref:inspect}{}UK AI Security Institute and collaborators. 2025. \textit{Inspect: a framework for large language model evaluations}. \href{https://inspect.aisi.org.uk/}{URL}.
    \item \hypertarget{ref:qwen}{}Qwen Team. 2024. \textit{Qwen2.5 Technical Report}. \href{https://arxiv.org/abs/2412.15115}{URL}.
    \item \hypertarget{ref:deepseek}{}DeepSeek-AI. 2026. \textit{DeepSeek-V4-Flash}. \href{https://huggingface.co/deepseek-ai/DeepSeek-V4-Flash}{URL}.
\end{enumerate}

\section*{LLM Usage Statement}

LLM-based tooling was used to accelerate repository scaffolding, dataset generation, judging passes, and manuscript editing. Experimental design, pre-registration, quality-gate definitions, analysis choices, and the adjudication of all results were human-led. Judge outputs were always dual-coded or validated against deterministic gates; no model was the sole validator of its own output.

\end{document}
